\documentclass[a4paper,11pt]{report}

% =====================================================
% Packages
% =====================================================
\usepackage[utf8]{inputenc}
\usepackage[T1]{fontenc}
\usepackage[french]{babel}
\usepackage{geometry}
\usepackage{listings}
\usepackage{xcolor}
\usepackage{hyperref}

\geometry{margin=2.5cm}

% =====================================================
% Configuration listings Java & Bash
% =====================================================
\definecolor{codegray}{rgb}{0.95,0.95,0.95}
\definecolor{darkgreen}{rgb}{0.0,0.5,0.0}
\definecolor{darkred}{rgb}{0.6,0.0,0.0}

\lstset{
  language=Java,
  backgroundcolor=\color{codegray},
  basicstyle=\ttfamily\small,
  keywordstyle=\color{blue},
  commentstyle=\color{darkgreen},
  stringstyle=\color{darkred},
  frame=single,
  breaklines=true,
  tabsize=2,
  showstringspaces=false
}

% =====================================================
% Titre
% =====================================================
\title{\textbf{Qualité de code en Java}\\
Bonnes pratiques, robustesse, tests et DevSecOps}
\author{Document pédagogique complet}

\begin{document}

\maketitle
\tableofcontents
\newpage

% =====================================================
\chapter{Introduction}

\section{Objectif du document}

Ce document a pour objectif de présenter les bonnes pratiques essentielles en langage Java à travers une approche comparative entre \textbf{mauvais} et \textbf{bon code}. Il vise à améliorer la qualité, la robustesse et la maintenabilité des applications en mettant en évidence les erreurs courantes et leurs corrections.

\section{Criticité}

Java offre un cadre plus sûr que des langages bas niveau, mais de mauvaises pratiques peuvent malgré tout conduire à :
\begin{itemize}
  \item des erreurs d'exécution comme les \texttt{NullPointerException}
  \item des fuites de ressources
  \item une dette technique importante
\end{itemize}

Adopter des règles strictes permet de produire un code lisible, fiable, testable et adapté à un contexte industriel.

\section{Contenu du document}

Ce document couvre les aspects fondamentaux de la qualité en Java :
\begin{itemize}
  \item Respect des conventions de nommage et de structure
  \item Gestion correcte de \texttt{null}, des exceptions et des ressources
  \item Bon usage des objets, collections et API modernes
  \item Écriture de code lisible, documenté et maintenable
\end{itemize}

\section{Approche pédagogique}

Chaque chapitre suit une structure simple et efficace :
\begin{itemize}
  \item un exemple de \textbf{mauvais code}
  \item une version corrigée (\textbf{bon code})
  \item une \textbf{explication} claire du problème et de la solution
\end{itemize}

Cette approche permet de comprendre rapidement les erreurs fréquentes et d'adopter les bons réflexes en situation réelle.

% =====================================================
\chapter{Conventions de nommage et casse}

\section{Mauvais code}
\begin{lstlisting}
package Com.MyProject.Utils;

class usermanager {
    int User_Age;
    void GetData() {}
}
\end{lstlisting}

\section{Bon code}
\begin{lstlisting}
package com.myproject.utils;

class UserManager {
    int userAge;

    void getData() {
    }
}
\end{lstlisting}

\section{Explication}
Les conventions Java standard imposent :
\begin{itemize}
  \item les packages en minuscules ;
  \item les classes en PascalCase ;
  \item les méthodes et variables en camelCase.
\end{itemize}

Ces règles rendent le code immédiatement lisible et cohérent entre développeurs.
Elles sont aussi vérifiées automatiquement par des outils comme Checkstyle.

% =====================================================
\chapter{Constantes et enums}

\section{Mauvais code}
\begin{lstlisting}
final int max = 7000;

enum state {
    on, off
}
\end{lstlisting}

\section{Bon code}
\begin{lstlisting}
static final int MAX_RPM = 7000;

enum State {
    ON,
    OFF
}
\end{lstlisting}

\section{Explication}
Les constantes doivent être facilement identifiables.
Une convention claire améliore la lisibilité et réduit les ambiguïtés entre variables ordinaires et invariants.

% =====================================================
\chapter{Commentaires utiles}

\section{Mauvais code}
\begin{lstlisting}
int i = 0; // initialise i a zero
\end{lstlisting}

\section{Bon code}
\begin{lstlisting}
int index = 0;
\end{lstlisting}

\section{Explication}
Un commentaire ne doit pas répéter ce que le code exprime déjà clairement.
Il doit expliquer pourquoi le code existe, quelle contrainte il respecte, ou quelle hypothèse il documente.

% =====================================================
\chapter{Commentaires orientés intention}

\section{Mauvais code}
\begin{lstlisting}
// wait a bit
Thread.sleep(200);
\end{lstlisting}

\section{Bon code}
\begin{lstlisting}
// Le capteur necessite 200 ms pour se stabiliser apres alimentation
Thread.sleep(200);
\end{lstlisting}

\section{Explication}
Un commentaire utile explique l'intention métier ou technique.
Il doit apporter une information que le code seul ne permet pas de déduire facilement.

% =====================================================
\chapter{Javadoc et contrat d'API}

\section{Mauvais code}
\begin{lstlisting}
/** get value */
int getValue() {
    return value;
}
\end{lstlisting}

\section{Bon code}
\begin{lstlisting}
/**
 * Retourne la tension mesuree par le capteur.
 *
 * @return tension en millivolts (>= 0)
 */
int getValue() {
    return value;
}
\end{lstlisting}

\section{Explication}
La Javadoc définit le contrat d'une méthode : unités, hypothèses, valeurs de retour, exceptions éventuelles.
Elle est essentielle pour les API publiques ou partagées.

% =====================================================
\chapter{Compilation stricte}

\section{Mauvais code}
\begin{lstlisting}
javac Main.java
\end{lstlisting}

\section{Bon code}
\begin{lstlisting}[language=bash]
javac -Xlint:all -Werror Main.java
\end{lstlisting}

\section{Explication}
Compiler avec des warnings activés permet de détecter plus tôt des problèmes de qualité.
Traiter les warnings comme des erreurs évite qu'ils s'accumulent dans le projet.

% =====================================================
\chapter{NullPointerException}

\section{Mauvais code}
\begin{lstlisting}
String role = getRole();
if (role.equals("ADMIN")) {
    grantAccess();
}
\end{lstlisting}

\section{Bon code}
\begin{lstlisting}
String role = getRole();
if ("ADMIN".equals(role)) {
    grantAccess();
}
\end{lstlisting}

\section{Explication}
Appeler une méthode sur une référence potentiellement nulle provoque une \texttt{NullPointerException}.
Placer le littéral à gauche évite ce risque.

% =====================================================
\chapter{Gestion des retours null avec Optional}

\section{Mauvais code}
\begin{lstlisting}
public User findUser(String id) {
    if (db.notContains(id)) {
        return null;
    }
    return db.get(id);
}
\end{lstlisting}

\section{Bon code}
\begin{lstlisting}
import java.util.Optional;

public Optional<User> findUser(String id) {
    if (db.notContains(id)) {
        return Optional.empty();
    }
    return Optional.of(db.get(id));
}
\end{lstlisting}

\section{Explication}
Retourner \texttt{null} oblige l'appelant à faire des vérifications manuelles, parfois oubliées.
\texttt{Optional} rend explicite la possibilité d'absence de valeur.

% =====================================================
\chapter{Comparaison d'objets}

\section{Mauvais code}
\begin{lstlisting}
String a = new String("x");
String b = new String("x");

if (a == b) {
}
\end{lstlisting}

\section{Bon code}
\begin{lstlisting}
if (a.equals(b)) {
}
\end{lstlisting}

\section{Explication}
L'opérateur \texttt{==} compare les références.
Pour comparer le contenu de deux objets, il faut utiliser \texttt{equals()}.

% =====================================================
\chapter{equals() et hashCode()}

\section{Mauvais code}
\begin{lstlisting}
class User {
    String email;

    public boolean equals(Object o) {
        return (o instanceof User u)
            && email.equals(u.email);
    }
}
\end{lstlisting}

\section{Bon code}
\begin{lstlisting}
class User {
    String email;

    public boolean equals(Object o) {
        return (o instanceof User u)
            && email.equals(u.email);
    }

    public int hashCode() {
        return email.hashCode();
    }
}
\end{lstlisting}

\section{Explication}
Si \texttt{equals()} est redéfini, \texttt{hashCode()} doit l'être aussi.
Sinon les collections de hachage comme \texttt{HashMap} ou \texttt{HashSet} se comportent mal.

% =====================================================
\chapter{Immuabilité et records}

\section{Mauvais code}
\begin{lstlisting}
public class Point {
    private final int x;
    private final int y;

    public Point(int x, int y) {
        this.x = x;
        this.y = y;
    }

    public int getX() { return x; }
    public int getY() { return y; }
}
\end{lstlisting}

\section{Bon code}
\begin{lstlisting}
public record Point(int x, int y) {}
\end{lstlisting}

\section{Explication}
Les \texttt{record} réduisent fortement le code répétitif.
Ils fournissent automatiquement un constructeur, des accesseurs, \texttt{equals()}, \texttt{hashCode()} et \texttt{toString()}.

% =====================================================
\chapter{Ressources non fermées}

\section{Mauvais code}
\begin{lstlisting}
FileInputStream in = new FileInputStream("data.txt");
byte[] data = in.readAllBytes();
in.close();
\end{lstlisting}

\section{Bon code}
\begin{lstlisting}
try (FileInputStream in = new FileInputStream("data.txt")) {
    byte[] data = in.readAllBytes();
}
\end{lstlisting}

\section{Explication}
Si une exception survient, une ressource ouverte manuellement peut ne jamais être libérée.
Le \texttt{try-with-resources} garantit la fermeture correcte.

% =====================================================
\chapter{Sortie standard et journalisation}

\section{Mauvais code}
\begin{lstlisting}
try {
    processData();
    System.out.println("Traitement termine");
} catch (Exception e) {
    e.printStackTrace();
}
\end{lstlisting}

\section{Bon code}
\begin{lstlisting}
import org.slf4j.Logger;
import org.slf4j.LoggerFactory;

private static final Logger LOGGER =
    LoggerFactory.getLogger(MyClass.class);

try {
    processData();
    LOGGER.info("Traitement termine avec succes");
} catch (Exception e) {
    LOGGER.error("Erreur lors du traitement des donnees", e);
}
\end{lstlisting}

\section{Explication}
\texttt{System.out.println()} et \texttt{printStackTrace()} sont peu adaptés à un code industriel.
Une bibliothèque de logging permet un meilleur contrôle, un meilleur formatage et une meilleure intégration aux outils de supervision.

% =====================================================
\chapter{Exceptions avalées}

\section{Mauvais code}
\begin{lstlisting}
try {
    save();
} catch (Exception e) {
}
\end{lstlisting}

\section{Bon code}
\begin{lstlisting}
try {
    save();
} catch (IOException e) {
    LOGGER.error("Erreur sauvegarde", e);
    throw new CustomException("Impossible de sauvegarder", e);
}
\end{lstlisting}

\section{Explication}
Un \texttt{catch} vide masque complètement l'erreur.
Il faut au minimum journaliser le problème, puis le propager ou le transformer en exception métier.

% =====================================================
\chapter{Concurrence : visibilité mémoire}

\section{Mauvais code}
\begin{lstlisting}
boolean stop = false;

void loop() {
    while (!stop) {}
}
\end{lstlisting}

\section{Bon code}
\begin{lstlisting}
volatile boolean stop = false;

void loop() {
    while (!stop) {}
}
\end{lstlisting}

\section{Explication}
Sans \texttt{volatile} ou synchronisation, un thread peut ne jamais observer la mise à jour faite par un autre thread.
Cela peut provoquer des boucles infinies.

% =====================================================
\chapter{Concurrence : opération non atomique}

\section{Mauvais code}
\begin{lstlisting}
int count = 0;

void inc() {
    count++;
}
\end{lstlisting}

\section{Bon code}
\begin{lstlisting}
import java.util.concurrent.atomic.AtomicInteger;

AtomicInteger count = new AtomicInteger(0);

void inc() {
    count.incrementAndGet();
}
\end{lstlisting}

\section{Explication}
L'opération \texttt{++} n'est pas atomique.
En environnement multithread, il faut utiliser des primitives adaptées comme \texttt{AtomicInteger}.

% =====================================================
\chapter{Modification de collection pendant itération}

\section{Mauvais code}
\begin{lstlisting}
for (String s : list) {
    if (s.isBlank()) {
        list.remove(s);
    }
}
\end{lstlisting}

\section{Bon code}
\begin{lstlisting}
list.removeIf(String::isBlank);
\end{lstlisting}

\section{Explication}
Modifier une collection pendant une boucle \texttt{for-each} provoque souvent une \texttt{ConcurrentModificationException}.
Il faut utiliser une méthode adaptée.

% =====================================================
\chapter{API Stream}

\section{Mauvais code}
\begin{lstlisting}
List<String> activeNames = new ArrayList<>();
for (User user : users) {
    if (user.isActive()) {
        activeNames.add(user.getName().toUpperCase());
    }
}
\end{lstlisting}

\section{Bon code}
\begin{lstlisting}
List<String> activeNames = users.stream()
    .filter(User::isActive)
    .map(user -> user.getName().toUpperCase())
    .toList();
\end{lstlisting}

\section{Explication}
L'API Stream permet d'exprimer plus directement l'intention : filtrer, transformer, collecter.
Elle réduit souvent le bruit technique des boucles manuelles.

% =====================================================
\chapter{Performance : String en boucle}

\section{Mauvais code}
\begin{lstlisting}
String s = "";
for (int i = 0; i < 10000; i++) {
    s += i;
}
\end{lstlisting}

\section{Bon code}
\begin{lstlisting}
StringBuilder sb = new StringBuilder();
for (int i = 0; i < 10000; i++) {
    sb.append(i);
}
String s = sb.toString();
\end{lstlisting}

\section{Explication}
Les objets \texttt{String} sont immuables.
Les concaténations répétées créent beaucoup d'objets temporaires ; \texttt{StringBuilder} est plus efficace.

% =====================================================
\chapter{Tests unitaires}

\section{Mauvais code}
\begin{lstlisting}
// Aucun test associe
int add(int a, int b) {
    return a + b;
}
\end{lstlisting}

\section{Bon code}
\begin{lstlisting}
import org.junit.jupiter.api.Test;
import static org.junit.jupiter.api.Assertions.assertEquals;

class CalculatorTest {
    @Test
    void shouldAddTwoNumbers() {
        assertEquals(4, add(2, 2));
    }
}
\end{lstlisting}

\section{Explication}
Les tests unitaires permettent de valider le comportement réel du code et de détecter les régressions.
Ils sont indispensables dans une chaîne de qualité moderne.

% =====================================================
\chapter{Couverture de code}

\section{Mauvais code}
\begin{lstlisting}
// Code de production jamais execute pendant les tests
\end{lstlisting}

\section{Bon code}
\begin{lstlisting}[language=bash]
mvn jacoco:report
\end{lstlisting}

\section{Explication}
La couverture mesure la part du code réellement exécutée pendant les tests.
Elle permet d'identifier des zones non testées.

% =====================================================
\chapter{Complexité excessive}

\section{Mauvais code}
\begin{lstlisting}
if (a) {
    if (b) {
        if (c) {
            if (d) {
                work();
            }
        }
    }
}
\end{lstlisting}

\section{Bon code}
\begin{lstlisting}
if (a && b && c && d) {
    work();
}
\end{lstlisting}

\section{Explication}
Une trop forte imbrication rend le code difficile à lire, à tester et à maintenir.
Réduire la complexité améliore la qualité globale du logiciel.

% =====================================================
\chapter{Analyse de style avec Checkstyle}

\section{Mauvais code}
\begin{lstlisting}
// Aucun controle automatique du style
\end{lstlisting}

\section{Bon code}
\begin{lstlisting}[language=bash]
mvn checkstyle:check
\end{lstlisting}

\section{Explication}
Checkstyle vérifie automatiquement le respect des conventions de style et de formatage.
Cela homogénéise le code à l'échelle du projet.

% =====================================================
\chapter{Analyse de bugs avec SpotBugs et PMD}

\section{Mauvais code}
\begin{lstlisting}
// Aucun scan de bugs ou mauvaises pratiques
\end{lstlisting}

\section{Bon code}
\begin{lstlisting}[language=bash]
mvn spotbugs:check
mvn pmd:check
\end{lstlisting}

\section{Explication}
Ces outils détectent des défauts fréquents, des mauvaises pratiques et certains problèmes de robustesse ou de sécurité.
Ils complètent utilement les tests.

% =====================================================
\chapter{Sécurité des dépendances}

\section{Mauvais code}
\begin{lstlisting}
// Dependances tierces jamais verifiees
\end{lstlisting}

\section{Bon code}
\begin{lstlisting}[language=bash]
mvn dependency-check:check
\end{lstlisting}

\section{Explication}
Une part importante des vulnérabilités vient des bibliothèques tierces.
Il faut analyser régulièrement les dépendances utilisées.

% =====================================================
\chapter{Inspection continue avec SonarQube}

\section{Mauvais code}
\begin{lstlisting}
// Aucun quality gate centralise
\end{lstlisting}

\section{Bon code}
\begin{lstlisting}[language=bash]
mvn clean verify sonar:sonar \
  -Dsonar.projectKey=mon-projet-java \
  -Dsonar.host.url=http://localhost:9000 \
  -Dsonar.token=mon_token_secret
\end{lstlisting}

\section{Explication}
SonarQube centralise la dette technique, les bugs, les vulnérabilités et la couverture.
Cela permet de mettre en place une vraie quality gate dans l'intégration continue.

% =====================================================
\chapter{Conclusion}

Un code Java de qualité professionnelle en approche DevSecOps repose sur :
\begin{itemize}
    \item des conventions de nommage et de style claires ;
    \item des API bien documentées ;
    \item une gestion maîtrisée de \texttt{null}, des exceptions et des ressources ;
    \item une attention particulière à la concurrence et à la lisibilité ;
    \item des tests unitaires et une couverture mesurée ;
    \item des analyses automatiques de style, de bugs, de sécurité et de dette technique.
\end{itemize}

Ces pratiques permettent de produire un code plus robuste, plus maintenable, plus sûr et plus industrialisable.

\end{document}